\subsection{Bài tập trắc nghiệm}
\ind{PHẦN I.} \inden{Câu trắc nghiệm nhiều phương án lựa chọn. Mỗi câu hỏi học sinh chỉ chọn một phương án.}\\
\setcounter{ex}{0}
\Opensolutionfile{ans}[ans/1D6-B1-1]
\begin{ex}%[2D2Y1]
	Cho $a$ là số thực dương. Đẳng thức nào sau đây đúng?
	\choice{$a^{x+y}=a^{x}+a^{y}$}
	{\True $\big( {a^x}\big)^y=a^{xy}$}
	{ $\big( {a^x}\big)^y=a^{x}.a^{y}$}
	{$a^{x-y}=a^{x}-a^{y}$}
\end{ex}

\begin{ex}%[2D2B1]
	Cho $a,b$ là các số thực dương khác $1$ và $x,y$ là các số thực. Đẳng thức nào sau đây đúng?
	\choice
	{\True $a^xa^y=a^{x+y}$}
	{$\dfrac{a^x}{a^y}=a^{\frac{x}{y}}$}
	{$a^xb^y=(ab)^{x+y}$}
	{$(a^x)^y=a^{x+y}$}
	\loigiai{
	}
\end{ex}
\begin{ex}
	Tìm số nhỏ hơn $1$ trong các số sau đây
	\choice
	{\True $\big(0,7\big)^{2017}$}
	{$\big(0,7\big)^{-2017}$}
	{$\big(1,7\big)^{2017}$}
	{$\big(2,7\big)^{2017}$}
\end{ex}

\begin{ex}
	Cho $\left(0{,}25\pi\right)^{\alpha}>\left(0{,}25\pi\right)^{\beta}$. Kết luận nào sau đây đúng?
	\choice
	{$\alpha\cdot \beta=1$}
	{$\alpha>\beta$}
	{$\alpha+\beta=0$}
	{\True $\alpha<\beta$}
	\loigiai{
		Do $0<0{,}25\pi<1$ nên $\left(0{,}25\pi\right)^{\alpha}>\left(0{,}25\pi\right)^{\beta}\Leftrightarrow \alpha <\beta$.
	}
\end{ex}

\begin{ex}
	Tính giá trị biểu thức $A = \dfrac{6^{3+\sqrt{5}}}{2^{2+\sqrt{5}}\cdot 3^{1+\sqrt{5}}}$.
	\choice
	{$1$}
	{$6^{-\sqrt{5}}$}
	{\True $18$}
	{$9$}
	\loigiai{
		Ta có $A = \dfrac{6^{3+\sqrt{5}}}{2^{2+\sqrt{5}}\cdot 3^{1+\sqrt{5}}}=\dfrac{2^{3+\sqrt{5}}\cdot 3^{3+\sqrt{5}}}{2^{2+\sqrt{5}}\cdot 3^{1+\sqrt{5}}}=2\cdot 3^2 =18$.
		%\shortsol{$18$}
	}
\end{ex}

\begin{ex}
	Giả sử $a$ là số thực dương. Biểu thức $\sqrt{a\sqrt[3]{a}}$ được viết dưới dạng $a^\alpha$. Khi đó giá trị $\alpha$ bằng bao nhiêu?
	\choice
	{\True $\alpha=\dfrac{2}{3}$}
	{$\alpha=\dfrac{11}{6}$}
	{$\alpha=\dfrac{1}{6}$}
	{$\alpha=\dfrac{5}{3}$}
	\loigiai
	{
		Ta có $\sqrt{a\sqrt[3]{a}}=\sqrt{\sqrt[3]{a^4}}=\sqrt[6]{a^4}=a^{\frac{2}{3}}$.
	}
\end{ex}

\begin{ex}
	Giả sử $x$ là số thực dương. Biểu thức $P=x \sqrt[5]x$ được viết dưới dạng lũy thừa với số mũ hữu tỉ là
	\boncot{$x^{\tfrac{11}{10}}$}{\True $x^{\tfrac{6}{5}}$}{$x^{\tfrac{1}{5}}$}{$x^{\tfrac{4}{5}}$}
\end{ex}


\begin{ex}
	Rút gọn biểu thức $P=x^{\tfrac{1}{3}}.\sqrt[6]{x}$ với $x>0$.
	\choice
	{$P=x^{\tfrac{1}{8}}$}
	{$P=x^2$}
	{\True $P=\sqrt{x}$}
	{$P=x^{\tfrac{2}{3}}$}
	\loigiai{Ta có: $P=x^{\tfrac{1}{3}}x^{\frac{1}{6}}=x^{\tfrac{1}{3}+\tfrac{1}{6}}=x^{\frac{1}{2}}=\sqrt{x}$.}
\end{ex}

\begin{ex}
	Rút gọn biểu thức $Q = b^{\frac{5}{3}} : \sqrt[3]{b}$, với $b > 0$.
	\choice
	{$Q = b^2$}
	{$Q = b^{\tfrac{5}{9}}$}
	{$Q = b^{-\tfrac{4}{3}}$}
	{\True $Q = b^{\tfrac{4}{3}}$}
	\loigiai {Ta có $Q = b^{\tfrac{5}{3}} : \sqrt[3]{b} = b^{\tfrac{5}{3}} : b^{\tfrac{1}{3}} = b^{\tfrac{5}{3} - \tfrac{1}{3}} = b^{\tfrac{4}{3}}$}
\end{ex}

\begin{ex}
	Rút gọn biểu thức $Q=\dfrac{b^{\frac{1}{3}}}{\sqrt[5]{b}}$, với $b>0$.
	\choice
	{$Q=b^{\frac{1}{15}}$}
	{$Q=b^{-\frac{2}{15}}$}
	{\True $Q=b^{\frac{2}{15}}$}
	{$Q=b^{\frac{5}{3}}$}
	\loigiai{$Q=b^{\frac{1}{3}-\frac{1}{5}}=b^{\frac{2}{15}}$
	}
\end{ex}

\begin{ex}
	Hãy viết biểu thức $L=\sqrt[3]{7.\sqrt[3]{7}}$ dưới dạng lũy thừa với số mũ hữu tỉ.
	\choice
	{$7^\frac{1}{2}$}
	{$7^{\frac{1}{18}}$}
	{\True $7^{\frac{4}{9}}$}
	{$7^{\frac{1}{27}}$}
\end{ex}

\begin{ex}
	Biến đổi $\sqrt[3]{x^5.\sqrt[4]{x}}, \left(x>0\right)$ thành dạng lũy thừa với số mũ hữu tỉ ta được
	\choice
	{${x}^{\tfrac{20}{3}}$}
	{${x}^{\tfrac{23}{12}}$}
	{\True ${x}^{\tfrac{21}{12}}$}
	{${x}^{\tfrac{12}{5}}$}
	\loigiai{Với $x>0$ ta có $\sqrt[3]{x^5.\sqrt[4]{x}}=\sqrt[3]{x^5\cdot x^{\tfrac{1}{4}}}=\sqrt[3]{x^{5+\tfrac{1}{4}}}=\left(x^{\tfrac{21}{4}}\right)^{\tfrac{1}{3}}=x^{\tfrac{21}{4}\cdot\tfrac{1}{3}}=x^{\tfrac{21}{12}}$.
	}
\end{ex}

\begin{ex}
	Viết biểu thức $A=\sqrt{a\sqrt{a\sqrt{a}}}:a^{\tfrac{11}{6}}$ $(a>0)$ dưới dạng số mũ lũy thừa hữu tỉ.
	\choice
	{\True $A=a^{-\tfrac{23}{24}}$}{$A=a^{\tfrac{21}{24}}$}{$A=a^{\tfrac{23}{24}}$}{$A=a^{-\tfrac{1}{12}}$}
\end{ex}

\begin{ex}
	Cho $a^{2b}=5$. Tính $2.a^{6b}$.
	\choice
	{$120$}
	{\True $250$}
	{$15$}
	{$125$}
	\loigiai{
		Xét $2.a^{6b}=2 \cdot (a^{2b})^3=2 \cdot 5^3 = 250$.}
\end{ex}

\begin{ex}
	Cho hai số dương $a$ và $b$ thỏa mãn $a^{\tfrac{1}{2}}=3$, $b^{\tfrac{1}{3}}=2$. Tính giá trị của tổng $S=a+b$.
	\choice
	{$5$}
	{$13$}
	{\True $17$}
	{$31$}
\end{ex}

\begin{ex}
	Giá trị của biểu thức $P=\left( 7+4\sqrt{3} \right)^{2017}\left( 7-4\sqrt{3} \right)^{2016}$ bằng
	\choice
	{$1$}
	{$7-4\sqrt{3}$}
	{\True $7+4\sqrt{3}$}
	{$\left( 7+4\sqrt{3} \right)^{2016}$}
	\loigiai{
		Ta có: 
		\allowdisplaybreaks
		\begin{eqnarray*}
			P & =& 	\left[ \left( 7+4\sqrt{3} \right)^{2016}\left( 7-4\sqrt{3} \right)^{2016} \right]\left( 7+4\sqrt{3} \right)\\
			&=& \left[ \left( 7+4\sqrt{3} \right)\left( 7-4\sqrt{3} \right) \right]^{2016}\left( 7+4\sqrt{3} \right)\\
			&=& (49-48)^{2016} \left( 7+4\sqrt{3} \right)=7+4\sqrt{3}.
		\end{eqnarray*}
	}
\end{ex}

\begin{ex}%[Bắc Đỗ - DA1]% %[2D2K1-1]%
	Giá trị của biểu thức $P=\left( 9+4\sqrt{5} \right)^{2017}\left( 9-4\sqrt{5} \right)^{2016}$ bằng
	\choice
	{\True $9+4\sqrt{5}$}
	{$1$}
	{$\left( 9-4\sqrt{5} \right)^{2016}$}
	{$9-4\sqrt{5}$}
	\loigiai{
		Ta có:
		\allowdisplaybreaks
		\begin{eqnarray*}
			P & =& \left[ \left( 9+4\sqrt{5} \right)^{2016}\left( 9-4\sqrt{5} \right)^{2016}\right] \left( 9+4\sqrt{5} \right)\\
			&=& \left[ \left( 9+4\sqrt{5} \right)\left( 9-4\sqrt{5} \right) \right]^{2016}\left( 9+4\sqrt{5} \right)\\
			&=& (81-80)^{2016} \left( 9+4\sqrt{5} \right)=9+4\sqrt{5}.
		\end{eqnarray*}
	}
\end{ex}

\begin{ex}%[Bắc Đỗ - DA1] %[2D2K1-1]
	Giá trị của biểu thức $P=\left( 1+\sqrt{3} \right)^{2016}\left( 3-\sqrt{3} \right)^{2016}$ bằng
	\choice
	{\True $12^{1008}$}
	{$4^{1008}$}
	{$\left( 1+\sqrt{3} \right)^{1008}$}
	{$\left( 3-\sqrt{3} \right)^{1008}$}
	\loigiai{
		Ta có: 
		\allowdisplaybreaks
		\begin{eqnarray*}
			P & =&\left[ \left( 1+\sqrt{3} \right)\left( 3-\sqrt{3} \right) \right]^{2016}\\
			&=& \left( 2 \sqrt{3}\right)^{2016}\\
			&=& 12^{1008}.
		\end{eqnarray*}
	}
\end{ex}


\begin{ex}
	Nếu $\big(a-2\big)^{-\tfrac{1}{4}} \le \big(a-2\big)^{-\tfrac{1}{3}}$ thì khẳng định nào sau đây đúng?
	\choice
	{$ a>3 $}
	{$ a<3 $}
	{\True $ 2<a<3 $}
	{$ a>2 $}
\end{ex}

\begin{ex}
	Cho $\big(a+1\big)^{-\tfrac{2}{3}} < \big(a+1\big)^{-\tfrac{1}{3}}$. Kết luận nào sau đây đúng?
	\choice
	{\True $a>0$}
	{$-1<a<0$}
	{$a \geq -1$}
	{$a \geq 0$}
\end{ex}

\Closesolutionfile{ans}

\ind{PHẦN II.} \inden{Câu trắc nghiệm đúng sai. Trong mỗi ý a), b), c), d) ở mỗi câu, học sinh chọn đúng hoặc sai.}\\
\Opensolutionfile{ans}[ans/1D6-B1-2]

\begin{ex}
	Cho ${a}$ là số thực dương tùy ý. Xét tính đúng sai của các khẳng định sau:
	\choiceTF
	{ \True $\sqrt[12] { a^{6} }= a^{\frac{1}{2}}$ }
	{ \True $(a^{4})^{11}=a^{44}$ }
	{ $a^{14}.a^{5}=a^{70}$ }
	{ $\dfrac{a^{2} } { a^{6} }=a^{\frac{1}{3}}$ }
	\loigiai{ 
		a) $\sqrt[12] {a^{6}}= a^{\frac{1}{2}}$ là khẳng định đúng.\\ 
		b) $(a^{4})^{11}=a^{44}$ là khẳng định đúng vì $(a^{4})^{11}=a^{4.11}=a^{44}$.\\ 
		c) $a^{14}.a^{5}=a^{70}$ là khẳng định sai vì $a^{14}.a^{5}=a^{14+5}=a^{19}$.\\ 
		d) $\dfrac{a^{2} } { a^{6} }=a^{\frac{1}{3}}$ là khẳng định sai vì $\dfrac{a^{2} } {a^{6} }=a^{2-6}=a^{-4}$.\\ 
		
}\end{ex}

\begin{ex}%[1D6H1-2]%::Cau 3::
	Cho $a=2$. Xét tính đúng sai của các khẳng định sau:
	\choiceTF
	{\True $a^3=8$}
	{\True $a^{-1}=\dfrac{1}{a}$}
	{$a^{\frac{1}{2}}\cdot a^{\frac{3}{2}}=a^3$}
	{\True Viết biểu thức $\dfrac{\sqrt{2\sqrt[3]{4}}}{16^{0{,}75}}$ về dạng lũy thừa $a^m$ ta được $m=-\dfrac{13}{6}$}
	\loigiai{
		\begin{itemchoice}
			\itemch \bf{Đúng}. Vì $a^3=2^3=8$.
			\itemch \bf{Đúng}. Vì $a=2\Rightarrow a\ne 0\Rightarrow a^{-1}=\dfrac{1}{a}$. (theo định nghĩa luỹ thừa với số mũ nguyên âm).
			\itemch \bf{Sai}. Vì $a^{\frac{1}{2}}\cdot a^{\frac{3}{2}}=a^{\frac{1}{2}+\frac{3}{2}}=a^{\frac{4}{2}}=a^2$.
			\itemch \bf{Đúng}. Vì $\dfrac{\sqrt{2\sqrt[3]{4}}}{16^{0{,}75}}
			=\dfrac{\sqrt{2\cdot 2^{\frac{2}{3}}}}{\left(2^4\right)^{\frac{3}{4}}}
			=\dfrac{\sqrt{2^{\frac{5}{3}}}}{2^3}=\dfrac{2^{\frac{5}{6}}}{2^3}
			=2^{\frac{5}{6}-3}=2^{\frac{-13}{6}}$.
		\end{itemchoice}
	}
\end{ex}

\begin{ex}
	Cho ${a,b}$ là các số thực dương tùy ý. Xét tính đúng sai của các khẳng định sau:
	\choiceTF
	{ \True $\sqrt[7]{\sqrt[10] { a} }=\sqrt[70]{a}$ }
	{ $\sqrt[3]{a}. \sqrt[3]{b}=\sqrt[3]{a + b}$ }
	{ \True $\sqrt[7]{a^{2}}= a^{\tfrac{2}{7}}$ }
	{ \True $\frac{\sqrt[4]{a}}{\sqrt[4]{b}}=\sqrt[4]{\frac{a}{b}}$ }
	\loigiai{ 
		a) $\sqrt[7]{\sqrt[10] { a} }=\sqrt[70]{a}$ là khẳng định đúng.\\ 
		b) $\sqrt[3]{a}. \sqrt[3]{b}=\sqrt[3]{a + b}$ là khẳng định sai vì $\sqrt[3]{a}. \sqrt[3]{b}=\sqrt[3]{a b}$.\\ 
		c) $\sqrt[7]a^{2}= a^{\tfrac{2}{7}}$ là khẳng định đúng.\\ 
		d) $\frac{\sqrt[4]{a}}{\sqrt[4]{b}}=\sqrt[4]{\frac{a}{b}}$ là khẳng định đúng.\\ 
		
}\end{ex}

\begin{ex}
	Cho các biều thức $A=\left(a^2\right)^{3+2 \sqrt{2}} \cdot a^{1-\sqrt{2}} \cdot a^{-4-\sqrt{2}}$ với $a>0$ và $B=\sqrt[4]{x \cdot \sqrt[3]{x^2 \cdot \sqrt{x^3}}}$ với $x>0$. Xét tính đúng sai của các khẳng định sau:
	\choiceTF
	{\True Với $a=2$ thì $A<57$}
	{Với $x=2$ thì $B>2$}
	{\True Khi $x=a$ thì $A \cdot B=a^{\frac{39+26 \sqrt{2}}{24}}$}
	{\True Khi $x=a$ thì $A: B=a^{\frac{72+48 \sqrt{2}}{13}}$}
	\loigiai{
		
	}
\end{ex}

\Closesolutionfile{ans}

\ind{PHẦN III.} \inden{Câu trắc nghiệm trả lời ngắn.}\\
\Opensolutionfile{ans}[ans/1D6-B1-3]

\begin{ex}%[1D4K1-2]
	Biết rằng $4^x=25^{y}=10$, với $x,y \in \mathbb{R}$. Tính giá trị của biểu thức $\dfrac{1}{x}+\dfrac{1}{y}$.\\
	\shortans[oly]{$2$}
	\loigiai{
		Ta có $4^x=10 \Rightarrow 10^{\tfrac{1}{x}}=4$; $25^y=10 \Rightarrow 10^{\tfrac{1}{y}}=25$.\\
		Suy ra $10^{\tfrac{1}{x}+\dfrac{1}{y}}=4\cdot25=100=10^2 \Rightarrow \dfrac{1}{x}+\dfrac{1}{y}=2$.	
	}
\end{ex}

\begin{ex}%[1D6H1-1]%::Cau 2::
	Trong năm $2\,022$, diện tích rừng trồng mới của tỉnh $A$ là $800$ ha. Giả sử diện tích rừng trồng mới của tỉnh $A$ mỗi năm tiếp theo đều tăng $7\%$ so với diện tích rừng trồng mới của năm liền trước. Hỏi năm $2\,030$, diện tích rừng trồng mới của tỉnh $A$ là bao nhiêu ha? (kết quả làm tròn đến hàng đơn vị)\\
	\shortans[oly]{$1\,375$}
	\loigiai{
		Năm $2\,030$ diện tích rừng trồng mới của tỉnh $A$ là $800(1+7\%)^8\approx 1\,375$ (ha).}
\end{ex}

\begin{ex}%[1D6H1-1]%::Cau 3::
	Năm $2\,020$ một hãng xe ô tô niêm yết giá bán loại xe $X$ là $850\,000\,000$ đồng và dự định trong năm tiếp theo, mỗi năm giảm $2\%$ giá bán của năm liền trước. Theo dự định đó, năm $2\,025$ hãng xe ô tô niêm yết giá bán xe $X$ là bao nhiêu triệu (kết quả làm tròn đến hàng đơn vị)?\\
	\shortans[oly]{$768$}
	\loigiai{
		Năm $2025$ giá bán loại xe $X$ là $850\,000\,000 (1-2\%)^5\approx 768$ (triệu đồng).}
\end{ex}

\begin{ex}%[1D4B1-4]
	Định luật thứ ba của Kepler về quỹ đạo chuyển động cho biết cách ước tính khoảng thời gian $P$ (tính theo năm Trái Đất) mà một hành tinh cần để hoàn thành một quỹ đạo quay quanh Mặt Trời. Khoảng thời gian đó được xác định bởi hàm số $P=d^{\tfrac{3}{2}}$, trong đó $d$ là khoảng cách từ hành tinh đó đến Mặt Trời tính theo đơn vị thiên văn $AU$ ($1\,AU$ là khoảng cách từ Trái Đất đến Mặt Trời, tức là $1\,AU$ khoảng $93\,000\,000$ dặm) (\textit{Nguồn$\colon$R.I.Charles et al.,Algebra $2$,Pearson}).\\
	Hỏi Sao Hỏa quay quanh Mặt Trời thì mất bao nhiêu năm Trái Đất (làm tròn kết quả đến hàng phần chục)? Biết khoảng cách từ Sao Hỏa đến Mặt Trời là $1{,}52$ AU.\\
	\shortans[oly]{$1,9$} 
	\loigiai{
		Thời gian để Sao Hỏa quay quanh Mặt Trời là
		$$P=d^{\tfrac{3}{2}}=1{,}52^{\tfrac{3}{2}}\approx1{,}9\text{ (năm Trái Đất)}.$$
	}	
\end{ex}


\begin{ex}
	Tại một xí nghiệp, công thức $P(t)=500 \cdot\left(\dfrac{1}{2}\right)^{\tfrac{t}{3}}$ được dùng để tính giá trị còn lại (tính theo triệu đồng) của một chiếc máy sau thời gian $t$ (tính theo năm) kể từ khi đưa vào sử dụng. Hỏi sau $1$ năm đưa vào sử dụng, giá trị còn lại của máy bằng bao nhiêu phần trăm so với ban đầu?(làm tròn kết quả đến hàng phần chục)\\
	\shortans[oly]{$79,4$} 
	\loigiai{
		Ban đầu giá trị của máy là $P_0=500\cdot \left(\dfrac{1}{2}\right)^0=500$.\\
		Giá trị còn lại của máy sau $1$ năm sử dụng: $P=500\cdot \left(\dfrac{1}{2}\right)^{\tfrac{1}{3}}=396{,}85$.\\
		Suy ra $\dfrac{P}{P_0}=79{,}4\%$.
	}	
\end{ex}

\begin{ex}%[1D6C1-2]%::Cau 6::
	Cho $f(x)=e^{\sqrt{1+\tfrac{1}{x^2}+\tfrac{1}{(x+1)^2}}}$. Biết rằng $f(1) \cdot f(2) \cdot f(3) \cdots f(2\,025)=e^{\frac{m}{n}}$ với $m$, $n$ là các số tự nhiên và $\dfrac{m}{n}$ là phân số tối giản. Tính $m-n^2$.\\
	\shortans[oly]{$-1$} 
	\loigiai{
		Đặt $g(x)=\sqrt{1+\dfrac{1}{x^2}+\dfrac{1}{(1+x)^2}}$.\\
		Với $x>0$, ta có\\
		$\begin{aligned}
			g(x) & =\sqrt{1+\dfrac{1}{x^2}+\dfrac{1}{(1+x)^2}}
			=\dfrac{\sqrt{x^2+(x+1)^2+x^2 (x+1)^2}}{x(x+1)}=\dfrac{\sqrt{(x^2+x+1)^2}}{x(x+1)} \\
			& =\dfrac{x^2+x+1}{x(x+1)}=1+\dfrac{1}{x(x+1)}=1+\dfrac{1}{x}-\dfrac{1}{x+1}.
		\end{aligned}$\\
		Suy ra 
		\begin{eqnarray*}
			&=& g(1)+g(2)+g(3)+\cdots+g(2\,025) \\
			&=&  \left(1+\dfrac{1}{1}-\dfrac{1}{2}\right)+\left(1+\dfrac{1}{2}-\dfrac{1}{3}\right)+\left(1+\dfrac{1}{3}-\dfrac{1}{4}\right)+\cdots+\left(1+\dfrac{1}{2\,025}-\dfrac{1}{2\,026}\right)\\
			&=& 2\,026-\dfrac{1}{2\,026}.
		\end{eqnarray*}
		Khi đó 
		\begin{eqnarray*}
			f(1) \cdot f(2) \cdot f(3) \cdots f(2\,025) &=& e^{g(1)+g(2)+g(3)+\cdots+g(2\,025)}\\
			&=& e^{2\,026-\frac{1}{2\,026}}\\
			&=& e^{\frac{2\,026^2-1}{2\,026}}=e^{\frac{m}{n}}.
		\end{eqnarray*}
		Do đó $m=2\,026^2-1$, $n=2\,026$.\\
		Vậy $m-n^2=2\,026^2-1-20\,268^2=-1$.
	}
\end{ex}

\Closesolutionfile{ans}